% GENERATED FILE. Edit canonical lessons, then run npm run generate:books.
% canonical-source-hash: fnv1a64:64985cd203b18ce0
% canonical-lessons: IT-C477-nostalgia, IT-C477-malinconia, IT-C477-rimpiangere, IT-C477-rimpianto, IT-C477-rimorso

\chapter{Ricordi e rimpianti}
\label{ch:it-ricordi-e-rimpianti}
\hlchaptermodality{477}

\begin{chapteropening}
\textbf{By the end of this chapter:} \emph{I can talk about missing the past and about regrets.}

Complete the last lesson of chapter 477: I can talk about missing the past and about regrets.
\end{chapteropening}

\section[la nostalgia]{la nostalgia — nostalgia, homesickness}

\label{lesson:IT-C477-nostalgia}



\begin{quote}
Before the new one: say the Italian for crying, then the Italian for to sob.
\end{quote}

\subsection*{You'll want to know: la nostalgia}
\textbf{la nostalgia} — \textquotedblleft{}nostalgia, homesickness\textquotedblright{}.

\begin{quote}
Quando vivevo all'estero avevo nostalgia di casa.
\end{quote}

\begin{quote}
\emph{When I lived abroad I was homesick.}
\end{quote}

\emph{Avere nostalgia di} — to miss (a place, a time).

\begin{cousinweb}
\emph{nostalgia} comes from New Latin \emph{nostalgia} (EtymDB), a modern coinage from Greek \emph{nóstos} \textquotedblleft{}return home\textquotedblright{} and \emph{álgos} \textquotedblleft{}pain\textquotedblright{}.
\end{cousinweb}

\begin{grammarlens}[title={What you've built}]
One more word for telling, explaining and describing at B1.
\end{grammarlens}

\subsection*{Your turn}
\begin{itemize}
\raggedright
  \item \emph{Say it:} \emph{la nostalgia}
  \item \emph{Say it:} \emph{la nostalgia}, once more
  \item \emph{Say it:} a sentence of your own with \emph{la nostalgia}
  \item \emph{Recall:} say the Italian for crying, then the Italian for to sob, then say \emph{la nostalgia} again
\end{itemize}

\begin{tcolorbox}[breakable,colback=teal!4,colframe=teal!35!black,title={Before you move on}]
What does la nostalgia mean? (\textquotedblleft{}Nostalgia, homesickness\textquotedblright{}.) Say it once more.
\end{tcolorbox}
\section[la malinconia]{la malinconia — melancholy, sadness}

\label{lesson:IT-C477-malinconia}



\begin{quote}
Before the new one: say the Italian for to sob, then the Italian for nostalgia.
\end{quote}

\subsection*{You'll want to know: la malinconia}
\textbf{la malinconia} — \textquotedblleft{}melancholy, sadness\textquotedblright{}.

\begin{quote}
La pioggia d'autunno mi mette malinconia.
\end{quote}

\begin{quote}
\emph{Autumn rain makes me melancholy.}
\end{quote}

The adjective is \emph{malinconico}.

\begin{cousinweb}
\emph{malinconia} comes from Late Latin \emph{melancholia} (EtymDB), from Greek \emph{melankholía} \textquotedblleft{}black bile\textquotedblright{}, which old medicine blamed for sadness.
\end{cousinweb}

\begin{grammarlens}[title={What you've built}]
One more word for telling, explaining and describing at B1.
\end{grammarlens}

\subsection*{Your turn}
\begin{itemize}
\raggedright
  \item \emph{Say it:} \emph{la malinconia}
  \item \emph{Say it:} \emph{la malinconia}, once more
  \item \emph{Say it:} a sentence of your own with \emph{la malinconia}
  \item \emph{Recall:} say the Italian for to sob, then the Italian for nostalgia, then say \emph{la malinconia} again
\end{itemize}

\begin{tcolorbox}[breakable,colback=teal!4,colframe=teal!35!black,title={Before you move on}]
What does la malinconia mean? (\textquotedblleft{}Melancholy, sadness\textquotedblright{}.) Say it once more.
\end{tcolorbox}
\section[rimpiangere]{rimpiangere — to miss, to regret (something lost)}

\label{lesson:IT-C477-rimpiangere}



\begin{quote}
Before the new one: say the Italian for nostalgia, then the Italian for melancholy.
\end{quote}

\subsection*{You'll want to know: rimpiangere}
\textbf{rimpiangere} — \textquotedblleft{}to miss, to regret (something lost)\textquotedblright{}.

\begin{quote}
Rimpiango gli anni dell'università.
\end{quote}

\begin{quote}
\emph{I miss my university years.}
\end{quote}

Built on \emph{piangere}; it takes a direct object: \emph{rimpiango quei giorni}.

\begin{cousinweb}
\emph{rimpiangere} is built from \emph{ri-}, \emph{in-} and \emph{piangere} \textquotedblleft{}to weep\textquotedblright{}, which is inherited from Latin \emph{plangere} \textquotedblleft{}to beat the breast in grief\textquotedblright{} (Lewis \& Short).
\end{cousinweb}

\begin{grammarlens}[title={What you've built}]
One more word for telling, explaining and describing at B1.
\end{grammarlens}

\subsection*{Your turn}
\begin{itemize}
\raggedright
  \item \emph{Say it:} \emph{rimpiangere}
  \item \emph{Say it:} \emph{rimpiangere}, once more
  \item \emph{Say it:} a sentence of your own with \emph{rimpiangere}
  \item \emph{Recall:} say the Italian for nostalgia, then the Italian for melancholy, then say \emph{rimpiangere} again
\end{itemize}

\begin{tcolorbox}[breakable,colback=teal!4,colframe=teal!35!black,title={Before you move on}]
What does rimpiangere mean? (\textquotedblleft{}To miss, to regret (something lost)\textquotedblright{}.) Say it once more.
\end{tcolorbox}
\section[il rimpianto]{il rimpianto — regret (for what is gone)}

\label{lesson:IT-C477-rimpianto}



\begin{quote}
Before the new one: say the Italian for melancholy, then the Italian for to miss.
\end{quote}

\subsection*{You'll want to know: il rimpianto}
\textbf{il rimpianto} — \textquotedblleft{}regret (for what is gone)\textquotedblright{}.

\begin{quote}
Non ho rimpianti: ho fatto le mie scelte.
\end{quote}

\begin{quote}
\emph{I have no regrets: I made my choices.}
\end{quote}

The past participle of \emph{rimpiangere}, turned into a noun.

\begin{cousinweb}
\emph{rimpianto} is the past participle of \emph{rimpiangere} \textquotedblleft{}to miss, to regret\textquotedblright{}, turned into a noun; the verb is built on \emph{piangere}, from Latin \emph{plangere} (Lewis \& Short).
\end{cousinweb}

\begin{grammarlens}[title={What you've built}]
One more word for telling, explaining and describing at B1.
\end{grammarlens}

\subsection*{Your turn}
\begin{itemize}
\raggedright
  \item \emph{Say it:} \emph{il rimpianto}
  \item \emph{Say it:} \emph{il rimpianto}, once more
  \item \emph{Say it:} a sentence of your own with \emph{il rimpianto}
  \item \emph{Recall:} say the Italian for melancholy, then the Italian for to miss, then say \emph{il rimpianto} again
\end{itemize}

\begin{tcolorbox}[breakable,colback=teal!4,colframe=teal!35!black,title={Before you move on}]
What does il rimpianto mean? (\textquotedblleft{}Regret (for what is gone)\textquotedblright{}.) Say it once more.
\end{tcolorbox}
\section[il rimorso]{il rimorso — remorse, guilt}

\label{lesson:IT-C477-rimorso}



\begin{quote}
Before the new one: say the Italian for to miss, then the Italian for regret.
\end{quote}

\subsection*{You'll want to know: il rimorso}
\textbf{il rimorso} — \textquotedblleft{}remorse, guilt\textquotedblright{}.

\begin{quote}
Ho il rimorso di non averla chiamata.
\end{quote}

\begin{quote}
\emph{I feel guilty for not calling her.}
\end{quote}

Remorse for something you did; for something lost, say \emph{il rimpianto}.

\begin{cousinweb}
\emph{rimorso} is the past participle of \emph{rimordere} \textquotedblleft{}to bite again\textquotedblright{}, from Latin \emph{remordēre} \textquotedblleft{}to vex, to torment\textquotedblright{} (Lewis \& Short), built on \emph{mordēre} \textquotedblleft{}to bite\textquotedblright{}: a conscience that bites.
\end{cousinweb}

\begin{grammarlens}[title={What you've built}]
One more word for telling, explaining and describing at B1.
\end{grammarlens}

\subsection*{Your turn}
\begin{itemize}
\raggedright
  \item \emph{Say it:} \emph{il rimorso}
  \item \emph{Say it:} \emph{il rimorso}, once more
  \item \emph{Say it:} a sentence of your own with \emph{il rimorso}
  \item \emph{Recall:} say the Italian for to miss, then the Italian for regret, then say \emph{il rimorso} again
\end{itemize}

\begin{tcolorbox}[breakable,colback=teal!4,colframe=teal!35!black,title={Before you move on}]
What does il rimorso mean? (\textquotedblleft{}Remorse, guilt\textquotedblright{}.) Say it once more.
\end{tcolorbox}
